\begin{resumo}[Abstract]
\begin{otherlanguage*}{english}
Patient safety and the retrospective tracking of adverse events in healthcare institutions represent essential and complex competencies in modern clinical education. The Global Trigger Tool (GTT) methodology, standardized by the Institute for Healthcare Improvement (IHI), mandates systematic audit of patient health records within a strict limit of twenty minutes per record, coupling trigger detection with root cause analysis via quality management tools. However, traditional approaches relying on physical paper or generic spreadsheets result in significant cognitive overload, time-induced anxiety, and visual fatigue for both students and faculty. In this scenario, SIGEA-GTT (Intelligent Adverse Event Management System -- Global Trigger Tool) was designed under the Human-Centred Design (HCD) paradigm, in compliance with the ISO 9241-210 standard. This document presents the formal specification of Personas, Empathy Maps, and User Journey Maps for the platform. The empirical design framework is grounded in Alan Cooper's Goal-Directed Design, Dave Gray and Alexander Osterwalder's Empathy Map Canvas, Donald Norman's Emotional Design, and Jakob Nielsen's Usability Engineering. Three archetypal personas are detailed: the primary persona \textit{Beatriz Matos} (Undergraduate Nursing Student in hospital internship), the secondary persona \textit{Prof. Dr. Ricardo Sobral} (Associate Professor and Senior Clinical Audit Specialist), and the tertiary persona \textit{Carlos Eduardo Santos} (Academic IT Systems Administrator), alongside a defined anti-persona (The Traditional Punitive Auditor). For each profile, six-quadrant empathy maps and native vector journey diagrams featuring emotional curves and touchpoint analyses are developed. These findings guide clinical ergonomic interface guidelines, desktop-oriented information density patterns, and four bidirectional traceability matrices mapping user pains and gains to the twenty-five functional requirements (RF01--RF25), ten non-functional requirements (RNF01--RNF10), and twenty-two use cases (UC01--UC22) of SIGEA-GTT. The monograph concludes with a formal heuristic evaluation and rigorous academic Quality Gate verification.

\textbf{Keywords}: Human-Centred Design. Personas. Empathy Map. User Journey. Global Trigger Tool. Software Ergonomics. Clinical Usability. SIGEA-GTT.
\end{otherlanguage*}
\end{resumo}
